\section{АНАЛИТИЧЕСКИЙ ОБЗОР}

\subsection{Особенности взаимодействия тренера и спортсмена}

Тренировочный процесс представляет собой повторяющийся цикл планирования, выполнения, отчётности и корректировки нагрузки. Тренер формирует задание, спортсмен выполняет его и передаёт обратную связь, после чего тренер оценивает состояние спортсмена и принимает решение о следующем этапе подготовки. Для эффективной работы важно, чтобы этот цикл был регулярным, прозрачным и основанным на достоверных данных.

В классической форме взаимодействие тренера и спортсмена может происходить очно: тренер объясняет задание на тренировке, наблюдает за выполнением и сразу корректирует действия спортсмена. Однако в реальных условиях такая модель не всегда достаточна. Спортсмены могут выполнять часть тренировок самостоятельно, находиться в разных местах, тренироваться в разное время или входить в разные группы. Поэтому тренеру необходим цифровой инструмент, позволяющий сопровождать тренировочный процесс вне очных занятий.

Цифровая поддержка тренировочного процесса должна решать несколько задач:

\begin{enumerate}
\def\labelenumi{\arabic{enumi}.}
\tightlist
\item
  передача тренировочных заданий спортсменам;
\item
  управление индивидуальными и групповыми назначениями;
\item
  получение отчётов о выполнении заданий;
\item
  хранение истории тренировок;
\item
  уведомление участников о важных событиях;
\item
  анализ накопленных данных;
\item
  разграничение ролей и прав доступа.
\end{enumerate}

Если рассматривать систему только как средство отправки сообщений, можно ограничиться мессенджером. Но при необходимости хранить историю тренировок, связывать отчёт с конкретным заданием и выполнять аналитику простого обмена сообщениями становится недостаточно. Требуется система, в которой данные имеют заранее определённую структуру.

\subsection{Использование мессенджеров}

Мессенджеры выступают наиболее распространённым инструментом цифрового взаимодействия благодаря своей массовости и интуитивно понятному интерфейсу, не требующему времени на освоение. Данный класс приложений позволяет тренеру оперативно создавать групповые чаты, рассылать тренировочные задания или консультировать спортсменов в индивидуальном порядке. В свою очередь, спортсмен получает возможность быстро предоставлять обратную связь, используя различные форматы данных.

Главное преимущество мессенджеров --- низкий порог входа. Для начала работы не требуется разрабатывать отдельную систему, создавать аккаунты в специализированном сервисе или обучать пользователей. Мессенджер удобен для оперативного общения и решения организационных вопросов.

Однако при использовании мессенджеров для ведения тренировочного процесса возникают существенные ограничения.

Ключевым ограничением является отсутствие стандартизированной структуры отчёта: спортсмен сам решает, что написать и в каком формате, в результате чего тренер получает разноформатные данные, непригодные для автоматической обработки.

Помимо этого, отчёты трудно связывать с конкретными заданиями. Если спортсмен прислал сообщение спустя несколько дней после нескольких назначений, тренер вынужден вручную восстанавливать соответствие, что особенно неудобно при групповой работе.

Существенной проблемой остаётся смешение тренировочных данных с общей перепиской: организационные сообщения, личные вопросы и отчёты оказываются в одном чате, что затрудняет поиск и систематизацию истории тренировок.

Наконец, мессенджеры не предоставляют аналитических инструментов: рассчитать процент выполнения заданий, динамику нагрузки или суммарную дистанцию без ручной обработки данных невозможно.

Следовательно, мессенджер эффективен как канал общения, но не является полноценной платформой для управления тренировочным процессом.

\subsection{Использование электронных таблиц}

Электронные таблицы частично решают проблему структуры данных. Тренер может создать таблицу с колонками для даты, задания, длительности, дистанции, пульса и комментария. Спортсмены могут заполнять строки после выполнения тренировок, а тренер --- просматривать данные и выполнять простые расчёты.

Преимущество таблиц заключается в том, что они позволяют хранить историю тренировок в более структурированном виде, чем мессенджеры. Таблицы также поддерживают фильтрацию, сортировку, формулы и визуализацию. Для небольших групп этого может быть достаточно.

Однако у таблиц также есть ограничения.

Ключевым ограничением является отсутствие полноценной ролевой модели: доступ можно настроить лишь на уровне документа или листа, но выразить предметные правила --- кто видит только свои задания, кто видит всю группу --- средствами таблицы затруднительно.

Вместе с тем таблицы не моделируют жизненный цикл задания: они хранят данные, но не отражают переходы между состояниями «назначено», «выполнено», «в архиве», что делает контроль за выполнением непрозрачным.

Дополнительную сложность создаёт отсутствие встроенного механизма уведомлений: тренер не получает системного оповещения об отправленном отчёте, а спортсмен может не заметить новое задание без дополнительного сообщения в мессенджере.

Наконец, при росте числа спортсменов таблица становится трудноуправляемой: возникает риск случайного изменения чужих данных, удаления строк и нарушения согласованности формата.

Таким образом, электронные таблицы лучше мессенджеров подходят для хранения истории, но не обеспечивают полноценную автоматизацию взаимодействия.

\subsection{Спортивные трекеры и приложения}

Спортивные трекеры и приложения предназначены для фиксации фактических показателей активности. Они могут автоматически собирать дистанцию, темп, пульс, маршрут, высоту, каденс и другие параметры. Для лёгкой атлетики такие инструменты полезны, потому что позволяют получить объективные данные о тренировке.

Сильной стороной спортивных трекеров является автоматизация измерений. Спортсмену не нужно вручную вводить все показатели, если устройство уже записало тренировку. Также многие приложения предоставляют графики и базовую аналитику.

Однако спортивные трекеры не полностью решают задачу взаимодействия тренера и спортсмена.

Ключевым ограничением является ориентированность трекеров на фиксацию факта активности, а не на управление назначениями: логика «задание назначено --- выполнено --- отчёт привязан к назначению» в большинстве таких сервисов отсутствует или реализована в ограниченном виде.

Помимо этого, трекеры, как правило, не учитывают организационный контекст тренера: тренировочные группы, индивидуальные планы, заявки на связь, ролевой доступ и уведомления о новых отчётах выходят за рамки их функциональности.

Существенным препятствием служит и закрытость большинства подобных сервисов: их API может быть ограничен, платным или недоступным для самостоятельного расширения, что исключает возможность интеграции в открытую платформу.

Наконец, не все виды тренировочной работы поддаются автоматической фиксации. Силовая подготовка, технические упражнения и восстановительные занятия требуют текстового комментария и субъективной оценки нагрузки, которые трекер не собирает.

Следовательно, спортивные трекеры полезны как потенциальный источник данных, но не заменяют backend-платформу взаимодействия тренера и спортсмена.

\subsection{Специализированные коммерческие платформы}

На рынке существуют специализированные платформы для тренеров, которые позволяют назначать тренировки, отслеживать выполнение, работать с группами и анализировать данные. К наиболее распространённым относятся TrainingPeaks, Final Surge и CoachNow. Такие решения ближе всего к требованиям предметной области.

Их преимуществами являются готовый пользовательский интерфейс, наличие функциональности для тренеров, интеграции с устройствами и развитая аналитика. Однако даже у таких платформ есть ряд принципиальных ограничений.

Ключевым из них является закрытость API и отсутствие возможности самостоятельной доработки: адаптировать платформу под специфические требования конкретного вида спорта или интегрировать её с внешними инструментами, как правило, невозможно.

Не менее значимой проблемой выступает зависимость от внешнего поставщика: пользователь не контролирует архитектуру, хранение данных, правила доступа и дальнейшее развитие продукта, что создаёт риски при изменении условий использования или прекращении работы сервиса.

Наконец, большинство профессиональных платформ ориентированы на коммерческое использование и недоступны для небольших групп или начинающих тренеров.

Таким образом, хотя функционал специализированных платформ доказывает востребованность подобных решений, их проприетарный характер не удовлетворяет требованиям расширяемости и независимости от поставщика, что обосновывает необходимость разработки собственной открытой системы.

\subsection{Сравнительный анализ подходов}

Сравнительный анализ проведён по критериям, важным для автоматизации тренировочного процесса. Результаты сравнения приведены в таблице~\ref{tab:approach-comparison}.

\begingroup
\TableFontPt{12}
\begin{longtable}{|
  >{\raggedright\arraybackslash}p{0.18\columnwidth}|
  >{\raggedright\arraybackslash}p{0.13\columnwidth}|
  >{\raggedright\arraybackslash}p{0.13\columnwidth}|
  >{\raggedright\arraybackslash}p{0.14\columnwidth}|
  >{\raggedright\arraybackslash}p{0.13\columnwidth}|
  >{\raggedright\arraybackslash}p{0.13\columnwidth}|}
\caption{Сравнение подходов к организации взаимодействия тренера и спортсмена}\label{tab:approach-comparison}\\
\hline
Критерий & Мессендж. & Эл.\ табл. & Трекеры & Коммерч.\linebreak платформы & CoachLink \\
\hline
\endfirsthead
\caption*{Продолжение таблицы \ref{tab:approach-comparison}}\\
\hline
Критерий & Мессендж. & Эл.\ табл. & Трекеры & Коммерч.\linebreak платформы & CoachLink \\
\hline
\endhead
\hline
\endlastfoot
Единый формат отчётов & Нет & Частично & Да, для активности & Да & Да \\ \hline
Связь отчёта с заданием & Нет & Частично & Ограниченно & Да & Да \\ \hline
Групповые назначения & Частично & Частично & Ограниченно & Да & Да \\ \hline
Роли тренера и спортсмена & Нет & Ограниченно & Частично & Да & Да \\ \hline
Уведомления о событиях & Да, но не профильные & Нет & Да & Да & Да \\ \hline
Хранение истории & Да, в разном формате & Да & Да & Да & Да \\ \hline
Аналитика & Нет & Базовая & Да, по активности & Да & Да \\ \hline
Открытость API & Нет & Ограниченно & Зависит от сервиса & Зависит от сервиса & Да \\ \hline
Возможность доработки & Низкая & Средняя & Низкая & Низк./сред. & Высокая \\ \hline
Самостоятельное\linebreak развёртывание & Нет & Частично & Нет & Обычно нет & Да \\ \hline
\end{longtable}
\endgroup

Анализ таблицы показывает, что ни один из распространённых подходов полностью не удовлетворяет требованиям открытой backend-платформы. Мессенджеры и таблицы удобны как вспомогательные инструменты, но не обеспечивают предметную модель. Спортивные трекеры полезны как источник данных, но не покрывают весь процесс взаимодействия. Специализированные коммерческие платформы функциональны, но ограничены с точки зрения открытости и самостоятельного развития.

CoachLink ориентируется на сочетание нескольких свойств: открытость, структурированность данных, ролевой доступ, групповые назначения, уведомления, аналитика и расширяемость. Это делает разработку собственной backend-платформы обоснованной.

\subsection{Функциональные требования}

На основе проведённого анализа предметной области можно сформулировать следующие функциональные требования к разрабатываемой системе.

\subsubsection{Пользователи и роли}

Подсистема управления пользователями должна обеспечивать:

\begin{itemize}
\tightlist
\item
  регистрацию и аутентификацию пользователей;
\item
  поддержку ролевой модели доступа с двумя базовыми ролями: «тренер» и «спортсмен»;
\item
  разграничение прав доступа к API-методам в зависимости от роли пользователя.
\end{itemize}

Такой подход позволяет изолировать функции, доступные разным категориям пользователей. Например, только тренер может создавать тренировочные планы и группы, а отправлять отчёт по назначению может только спортсмен, которому это назначение принадлежит.

\subsubsection{Связь тренер-спортсмен}

Подсистема управления связями между тренером и спортсменом должна обеспечивать:

\begin{itemize}
\tightlist
\item
  отправку спортсменом заявки на установление связи с тренером;
\item
  возможность принятия или отклонения входящей заявки тренером;
\item
  ограничение, при котором один спортсмен может быть связан не более чем с одним тренером одновременно.
\end{itemize}

Механизм подтверждаемых заявок защищает тренера от автоматического добавления спортсменов и позволяет явно контролировать состав подопечных. Ограничение связи в первой версии системы упрощает логику доступа и однозначно определяет тренера, ответственного за тренировочный процесс спортсмена.

\subsubsection{Тренировочные группы}

Подсистема управления тренировочными группами должна обеспечивать:

\begin{itemize}
\tightlist
\item
  создание и редактирование тренером тренировочных групп;
\item
  добавление спортсменов в группу только при наличии установленной связи с данным тренером.
\end{itemize}

Объединение спортсменов в группы необходимо для массового назначения одинаковых тренировочных планов нескольким спортсменам одновременно.

\subsubsection{Тренировочные планы и назначения}

Подсистема планирования тренировок должна обеспечивать:

\begin{itemize}
\tightlist
\item
  создание тренером тренировочных планов;
\item
  назначение тренировочных планов отдельным спортсменам или тренировочным группам;
\item
  просмотр списка назначений;
\item
  поддержку статусов назначений, позволяющих отличать активные, выполненные и архивные задания;
\item
  сохранение описаний тренировок как шаблонов для повторного использования.
\end{itemize}

Данный функционал автоматизирует основной цикл предметной области: постановку тренировочного задания, его выполнение спортсменом и последующий контроль со стороны тренера.

\subsubsection{Отчёты}

Подсистема отчётности должна обеспечивать:

\begin{itemize}
\tightlist
\item
  отправку спортсменом структурированного отчёта по конкретному тренировочному назначению;
\item
  хранение отчёта в едином формате;
\item
  поддержку полей отчёта, необходимых для первой версии системы: текстовый комментарий, длительность, субъективная нагрузка, дистанция и показатели пульса.
\end{itemize}

Единый формат отчётов устраняет разрозненность данных и делает собранную информацию пригодной для последующей автоматической обработки.

\subsubsection{Уведомления}

Подсистема уведомлений должна обеспечивать:

\begin{itemize}
\tightlist
\item
  создание уведомлений о важных событиях: заявке на связь, принятии заявки, назначении тренировки, удалении задания, отправке отчёта, добавлении или удалении из группы;
\item
  хранение уведомлений в базе данных;
\item
  предоставление уведомлений пользователям через API.
\end{itemize}

Система уведомлений обеспечивает своевременное информирование участников о ключевых событиях без необходимости постоянного мониторинга приложения.

\subsubsection{Аналитика}

Подсистема аналитики должна обеспечивать:

\begin{itemize}
\tightlist
\item
  агрегацию данных из отчётов спортсмена;
\item
  расчёт сводных показателей за выбранный период: количества отчётов, суммарной длительности, суммарной дистанции, средней субъективной нагрузки и средних показателей пульса;
\item
  вычисление процента выполненных назначений.
\end{itemize}

Агрегированные показатели позволяют тренеру оценивать динамику нагрузки и эффективность тренировочного процесса на основе накопленных структурированных данных.

\subsubsection{AI-рекомендации}

Подсистема интеллектуальной поддержки должна обеспечивать генерацию текстовой рекомендации по конкретному спортсмену на основе последних отчётов и сводной статистики.

AI-модуль рассматривается как вспомогательный инструмент поддержки принятия решений и не заменяет экспертную оценку тренера.

\subsection{Нефункциональные требования}

Помимо функциональных возможностей, backend-платформа должна удовлетворять ряду нефункциональных требований, связанных с архитектурой, сопровождением и эксплуатацией системы. К ним относятся:

\begin{enumerate}
\def\labelenumi{\arabic{enumi}.}
\tightlist
\item
  модульность архитектуры и возможность независимого развития основных компонентов системы;
\item
  наличие единого программного интерфейса для взаимодействия внешних клиентов с backend-платформой;
\item
  документирование публичных API-контрактов в машиночитаемом формате;
\item
  логическое разделение данных между функциональными областями системы;
\item
  надёжное хранение структурированных данных с поддержкой транзакционности;
\item
  поддержка асинхронной обработки событий между компонентами системы;
\item
  воспроизводимое локальное развёртывание системы;
\item
  возможность расширения предметной модели под другие виды спорта;
\item
  тестируемость бизнес-логики;
\item
  минимальная зависимость от конкретного клиентского приложения.
\end{enumerate}

\subsection{Выводы по главе 1}

В первой главе были рассмотрены основные способы организации цифрового взаимодействия тренера и спортсмена: мессенджеры, электронные таблицы, спортивные трекеры и специализированные коммерческие платформы. Анализ показал, что каждый из подходов решает только часть задачи.

Мессенджеры удобны для общения, но не обеспечивают структурированное хранение данных. Электронные таблицы позволяют хранить историю, но не обладают полноценной предметной логикой, уведомлениями и ролевой моделью. Спортивные трекеры собирают данные активности, но не организуют полный цикл назначения и отчётности. Специализированные коммерческие платформы функциональны, но ограничены с точки зрения открытости, самостоятельного развёртывания и доработки.

На основании анализа сформулированы требования к backend-платформе CoachLink. Система должна поддерживать роли пользователей, связи тренера и спортсмена, тренировочные группы, планы, назначения, стандартизированные отчёты, уведомления, аналитику и AI-рекомендации. Для удовлетворения этих требований в следующей главе предлагается микросервисная архитектура backend-платформы.
